\begin{table}[ht]
\centering
\footnotesize
\setlength{\tabcolsep}{4pt}
\caption{Total sample size ($N$) of the conventional iterative method ($N_{\rm conv}$) and the proposed closed-form formula ($N_{\rm prop}$), and the achieved co-primary power at $N_{\rm prop}$, for target powers $1 - \beta$ of $0.80$, $0.85$, and $0.90$. The significance level is $\alpha = 0.025$ (one-sided).}
\label{tab:ss_comparison}
\begin{threeparttable}
\begin{tabular}{ccccrrrrrrrrr}
\hline
 & & & & \multicolumn{3}{c}{$1 - \beta = 0.80$} & \multicolumn{3}{c}{$1 - \beta = 0.85$} & \multicolumn{3}{c}{$1 - \beta = 0.90$} \\
\cline{5-7} \cline{8-10} \cline{11-13}
$(\delta_1, \delta_2)$ & $(\sigma_1, \sigma_2)$ & $r$ & $\rho$ & $N_{\rm conv}$ & $N_{\rm prop}$ & power & $N_{\rm conv}$ & $N_{\rm prop}$ & power & $N_{\rm conv}$ & $N_{\rm prop}$ & power \\
\hline
(0.3, 0.3) & (1.0, 1.0) & 1 & 0.0 & 460 & 460 & 0.8022 & 508 & 508 & 0.8507 & 574 & 574 & 0.9003 \\
 &  &  & 0.3 & 448 & 448 & 0.8016 & 498 & 498 & 0.8506 & 566 & 566 & 0.9003 \\
 &  &  & 0.5 & 436 & 436 & 0.8011 & 488 & 488 & 0.8512 & 558 & 556 & 0.9000$^{\dagger}$ \\
 &  &  & 0.8 & 408 & 408 & 0.8010 & 460 & 460 & 0.8505 & 530 & 530 & 0.9000 \\
(0.3, 0.3) & (1.0, 1.0) & 2 & 0.0 & 516 & 516 & 0.8007 & 573 & 573 & 0.8519 & 648 & 648 & 0.9016 \\
 &  &  & 0.3 & 504 & 504 & 0.8016 & 561 & 561 & 0.8512 & 639 & 639 & 0.9015 \\
 &  &  & 0.5 & 492 & 492 & 0.8026 & 549 & 549 & 0.8512 & 627 & 627 & 0.9008 \\
 &  &  & 0.8 & 459 & 459 & 0.8010 & 519 & 519 & 0.8516 & 597 & 597 & 0.9004 \\
(0.3, 0.3) & (1.0, 1.5) & 1 & 0.0 & 804 & 804 & 0.8005 & 910 & 910 & 0.8503 & 1058 & 1058 & 0.9004 \\
 &  &  & 0.3 & 798 & 798 & 0.8004 & 906 & 906 & 0.8501 & 1056 & 1056 & 0.9003 \\
 &  &  & 0.5 & 794 & 794 & 0.8009 & 904 & 904 & 0.8506 & 1054 & 1054 & 0.9002 \\
 &  &  & 0.8 & 786 & 786 & 0.8002 & 900 & 900 & 0.8507 & 1052 & 1052 & 0.9003 \\
(0.3, 0.3) & (1.0, 1.5) & 2 & 0.0 & 906 & 906 & 0.8012 & 1026 & 1026 & 0.8511 & 1191 & 1191 & 0.9005 \\
 &  &  & 0.3 & 897 & 897 & 0.8000 & 1020 & 1020 & 0.8503 & 1188 & 1188 & 0.9003 \\
 &  &  & 0.5 & 894 & 894 & 0.8013 & 1017 & 1017 & 0.8506 & 1185 & 1185 & 0.9000 \\
 &  &  & 0.8 & 885 & 885 & 0.8006 & 1011 & 1011 & 0.8502 & 1185 & 1185 & 0.9007 \\
(0.7, 0.5) & (1.0, 1.0) & 1 & 0.0 & 132 & 132 & 0.8031 & 148 & 148 & 0.8510 & 172 & 172 & 0.9025 \\
 &  &  & 0.3 & 130 & 130 & 0.8013 & 148 & 148 & 0.8536 & 170 & 170 & 0.9000 \\
 &  &  & 0.5 & 130 & 130 & 0.8057 & 146 & 146 & 0.8511 & 170 & 170 & 0.9012 \\
 &  &  & 0.8 & 128 & 128 & 0.8062 & 144 & 144 & 0.8502 & 170 & 170 & 0.9029 \\
(0.7, 0.5) & (1.0, 1.0) & 2 & 0.0 & 150 & 150 & 0.8076 & 168 & 168 & 0.8544 & 192 & 192 & 0.9002 \\
 &  &  & 0.3 & 147 & 147 & 0.8035 & 165 & 165 & 0.8502 & 192 & 192 & 0.9012 \\
 &  &  & 0.5 & 147 & 147 & 0.8078 & 165 & 165 & 0.8527 & 192 & 192 & 0.9023 \\
 &  &  & 0.8 & 144 & 144 & 0.8062 & 162 & 162 & 0.8502 & 192 & 192 & 0.9040 \\
(0.7, 0.5) & (1.0, 1.5) & 1 & 0.0 & 284 & 284 & 0.8020 & 324 & 324 & 0.8508 & 380 & 380 & 0.9013 \\
 &  &  & 0.3 & 284 & 284 & 0.8020 & 324 & 324 & 0.8508 & 380 & 380 & 0.9013 \\
 &  &  & 0.5 & 284 & 284 & 0.8020 & 324 & 324 & 0.8508 & 380 & 380 & 0.9013 \\
 &  &  & 0.8 & 284 & 284 & 0.8020 & 324 & 324 & 0.8508 & 380 & 380 & 0.9013 \\
(0.7, 0.5) & (1.0, 1.5) & 2 & 0.0 & 318 & 318 & 0.8001 & 366 & 366 & 0.8523 & 426 & 426 & 0.9003 \\
 &  &  & 0.3 & 318 & 318 & 0.8001 & 366 & 366 & 0.8523 & 426 & 426 & 0.9003 \\
 &  &  & 0.5 & 318 & 318 & 0.8001 & 366 & 366 & 0.8523 & 426 & 426 & 0.9003 \\
 &  &  & 0.8 & 318 & 318 & 0.8001 & 366 & 366 & 0.8523 & 426 & 426 & 0.9003 \\
\hline
\end{tabular}
\begin{tablenotes}
\footnotesize
\item[] Achieved power is the bivariate normal power at the group sizes $(n_1, n_2)$ of the proposed formula, computed with the \texttt{pbivnorm} package. The proposed formula uses the 5-point Gauss--Legendre rule.
\item[$\dagger$] Achieved power below the target: at $(\delta_1, \delta_2, \sigma_1, \sigma_2, r, \rho) = (0.3, 0.3, 1.0, 1.0, 1, 0.5)$ and $1 - \beta = 0.90$, the achieved power at $N_{\rm prop} = 556$ is 0.8999989, whereas $N_{\rm conv} = 558$.
\end{tablenotes}
\end{threeparttable}
\end{table}
